% set up
\begin{center}
    {\fontsize{18pt}{0pt} \MakeUppercase{\textbf{abstract}}}
\vspace*{24pt}
\end{center}
\addcontentsline{toc}{chapter}{Abstract}
%-------------------------------------------------------

Collectors of trading cards need to identify the exact card in a photo: the same Pok\'emon appears in many versions, sets reprint identical artwork, and phone photos add perspective, glare and blur. This report describes a pipeline built for this task, its systematic evaluation, and an application that puts it into practice. A classical detector finds the card outline (Canny edges, contours, convex hull and quality checks) and warps the card to a canonical $500\times700$ view with a homography. The rectified card is then matched against a reference database of \res{db-cards} official scans from \res{db-sets} recent sets by interchangeable retrieval methods: a pixel-thumbnail baseline, SIFT and ORB keypoints verified with RANSAC, frozen ImageNet ResNet-50 embeddings, frozen CLIP ViT-B/32 embeddings, and a hybrid that re-ranks an embedding shortlist with SIFT and RANSAC. Reprints are handled with artwork groups found by perceptual hashing (\res{db-artgroups} distinct artworks), and every result is reported both strictly (exact card) and leniently (same artwork), with 95\% bootstrap intervals.

On synthetic single-card photos with known ground truth and five distortion levels, exhaustive keypoint matching was the most accurate (strict top-1 at the medium level: ORB \res{orb-s3-top1}, SIFT \res{sift-s3-top1}), and its errors at mild distortion were only reprints with identical artwork. Re-ranking a 20-card CLIP shortlist with SIFT raised CLIP from \res{clip-s3-top1} to \res{hybrid_clip_k20-s3-top1}, matching exhaustive SIFT at about \res{hybrid_clip_k20-s3-ms}~ms instead of \res{sift-s3-ms}~ms per photo, whereas ResNet-50 alone reached only \res{resnet-s3-top1}. Without rectification the global methods failed almost completely and SIFT dropped to \res{norect-sift-top1}; under strong distortion or with overlapping cards most errors came from not finding the card outline. In a separate test with English, Spanish and Japanese cards photographed at phone resolution, SIFT with RANSAC was nearly language-independent (lenient top-1 between \res{lang-sift-ja-hr-len} and \res{lang-sift-en-hr-len}), while CLIP alone dropped to \res{lang-clip-ja-hr-len} on Japanese cards; re-ranking the CLIP shortlist with SIFT gave \res{lang-hybrid_clip+rot_k20-ja-hr-len} with a median of about \res{lang-hybrid_clip+rot_k20-ms}~ms per photo, compared with \res{lang-sift-ms}~ms for exhaustive SIFT. The pipeline runs in a local web application, which also recognises the \res{db-ja-cards} cards of a Japanese set and was used to catalogue \res{app-n-items} real cards. On their saved phone photos, SIFT, ORB, CLIP and the CLIP hybrid identified every card, ResNet-50 \res{app-resnet-hits} and the pixel baseline only \res{app-pixel-hits}, which suggests that the synthetic benchmark favours pixel-level similarity. Two extensions are also evaluated: printed-language recognition with OCR (accuracy \res{langclf-v2-test-acc} on the test split) and a visual condition estimate, which works on synthetically damaged cards but not yet on real photos of sleeved cards.
